\bookchapter{Chapter 29: Digital Wellbeing — Building a Healthy Relationship with Your Phone}

\chapterintro{Have you ever experienced this: you only meant to scroll short videos for 5 minutes, then looked up and realized two hours had passed? Or lying in bed at night scrolling your phone, getting more and more awake, and finally falling asleep at 2 AM? Phones have changed our lives, but they're also quietly consuming our time and health. This chapter isn't about “quitting phones” — it's about helping you take back control.}

\booksection{I. Check Your “Screen Time” — Numbers Are a Wake-Up Call}

Before making changes, understand the current situation. Your phone has a feature many people have never opened called “Screen Time” (iPhone) or “Digital Wellbeing” (Android). It tells you exactly how much time you spend on your phone each day and which apps that time goes to.

\begin{itemize}
  \item \textbf{iPhone}: Settings > Screen Time > See All Activity. You'll see daily unlock counts, pickup counts, and time spent in each app.
  \item \textbf{Android}: Settings > Digital Wellbeing \& Parental Controls. Huawei calls it "Healthy Phone Usage," Xiaomi calls it "Screen Time Management." Similar features, all showing daily usage data.
\end{itemize}

\textbf{80\% of people who see this data for the first time are shocked}. A friend of mine discovered after checking that they spent over 4 hours a day on Douyin — that's 28 hours a week, equivalent to a part-time job. This kind of “digital shock” is the first step toward change.

\booksection{II. Set App Time Limits — Install a Speed Limiter for Yourself}

Once you know where the problem is, you can set targeted limits. No need to delete apps outright — just set a “daily allowance”:

\begin{enumerate}
  \item \textbf{iPhone}: Settings > Screen Time > App Limits > Add Limit. Select the app you want to restrict (e.g., Douyin) and set 30 minutes per day. When the time's up, the app icon grays out, and opening it requires an additional password entry.
  \item \textbf{Android}: Settings > Digital Wellbeing > App Timers. Select an app and set a daily usage limit; when time's up, the app icon grays out.
  \item \textbf{Tip}: Set the limit password to something you don't use often and is hard to remember. If it's too easy to enter, you'll be tempted to "extend" every time, making the limit useless.
\end{enumerate}

\booksection{III. The Right Way to Use “Do Not Disturb” — Not Muting Notifications, But Muting Anxiety}

Many people's phones buzz with notifications dozens of times a day, and each buzz interrupts their attention. You think you're just “glancing at a notification,” but it actually takes 15-20 minutes to refocus after each interruption. Do Not Disturb isn't about “going offline” — it's about helping you regain control of your attention.

\begin{itemize}
  \item \textbf{Schedule It}: Set Do Not Disturb to automatically turn on from 10 PM to 7 AM. This way, your sleep won't be interrupted by a spam text message.
  \item \textbf{Allow Important Contacts Through}: In Do Not Disturb settings, set "Allow Calls From" to "Favorites" or enable "Repeated Calls" (a second call within 3 minutes will ring normally). This way you won't miss emergency calls from family, but won't be bothered by sales calls.
  \item \textbf{Work Focus Mode}: iPhone's "Focus" mode lets you set custom scenarios — for example, a "Work" focus only allows work-related app notifications, while a "Reading" focus silences all notifications.
\end{itemize}

\booksection{IV. Eye Protection Relies on Habits, Not Apps — The 20-20-20 Rule}

Prolonged screen time inevitably causes dry, sore eyes. Eye protection mode (blue light filter) helps somewhat, but what really works is changing your eye usage habits:

\begin{tcbinfo}{The 20-20-20 Rule}
Every\textbf{20 minutes}of screen time, look at something\textbf{20 feet (about 6 meters)}away for\textbf{20 seconds}. This simple action effectively relieves ciliary muscle fatigue.

\end{tcbinfo}

\begin{itemize}
  \item \textbf{Enable Eye Protection Mode}: Settings > Display > Eye Protection (or Night Shift). When enabled, the screen's color temperature warms up, reducing blue light. Recommended to set it to auto-enable from sunset to sunrise.
  \item \textbf{Adjust Screen Brightness}: When the environment is dark, the screen should be dim too. Many people keep their screen too bright in the dark — this is the most damaging behavior for your eyes. Enable "Auto-Brightness" and make sure it works properly.
  \item \textbf{Maintain Distance}: Keep the phone at least 30 cm from your eyes. Many people hold their phones very close, which keeps the ciliary muscle in a state of tension, accelerating nearsightedness.
\end{itemize}

\booksection{V. Neck Rescue Guide for “Phone Benders”}

When looking down at your phone, your neck bears 5-6 times the pressure of normal standing. Long-term neck bending can straighten the cervical curve, cause bone spurs, dizziness, and hand numbness. Here are some simple improvement methods:

\begin{itemize}
  \item \textbf{Raise Your Phone}: When looking at your phone, try to hold it at eye level rather than looking down at your lap. Your arms may get tired, but your neck will thank you.
  \item \textbf{Move Your Neck Every 30 Minutes}: Slowly turn your head left to the limit, hold for 3 seconds, return to center; then turn right, hold for 3 seconds. Repeat 5 times. Then do the "rice character exercise" — trace the character 米 ("rice") in the air with your chin.
  \item \textbf{Wall Standing for 3 Minutes}: Press the back of your head, shoulder blades, buttocks, and heels against a wall. Stand for 3 minutes daily to effectively correct the hunched, head-down posture.
\end{itemize}

\booksection{VI. Sleep and Phone — The Bedroom Habit Most Worth Changing}

Blue light from your phone screen suppresses melatonin production, making you “more awake the more you look.” Many people feel “I'll get sleepy as I scroll,” but they're actually being tricked by blue light — when they finally put the phone down, they can't fall asleep. Here are improvement suggestions:

\begin{itemize}
  \item \textbf{No Phone for 1 Hour Before Bed}: This is the most effective yet hardest to follow. If you can't manage it, at least turn on "Night Shift/Eye Protection Mode" to minimize blue light.
  \item \textbf{Put the Phone Out of Reach}: Put the charger at the other end of the bedroom, not on the nightstand. This way you won't be tempted to pick it up and scroll, and when the alarm goes off, you have to get up to turn it off — killing two birds with one stone.
  \item \textbf{Enable "Sleep Mode"}: iPhone's "Sleep" Focus can automatically lock the screen and simplify the lock screen during sleep hours; Android brands have similar features.
  \item \textbf{Use an Alarm Clock Instead of Your Phone}: Buy a cheap physical alarm clock for a dozen or so yuan. Stop using your phone as an alarm, and you'll escape the trap of "just checking the time, then can't help scrolling."
\end{itemize}

\booksection{VII. Information Anxiety — Learn “Selective Blindness”}

The amount of information we receive daily equals what people in ancient times received in a lifetime. But the human brain wasn't designed to process this much information — overload leads to anxiety, scattered attention, and decision fatigue. The solution isn't “getting more information” but “actively filtering”:

\begin{itemize}
  \item \textbf{Unfollow}: Open your list of followed accounts, bloggers, and creators, and ask yourself one question: "Did I get any real value from this in the past month?" If not, unfollow. Quality over quantity for information sources.
  \item \textbf{Turn Off Non-Essential App Notifications}: Turn off notifications for all apps except WeChat, phone, SMS, and calendar. News doesn't need push notifications — you'll see it when you actively look.
  \item \textbf{Set an "Information Fast Day"}: Pick one day a week (e.g., Sunday) to completely avoid phone news and short videos. You'll find the sky doesn't fall, and you'll feel noticeably better on Monday.
\end{itemize}

\begin{tcbpractice}{Hands-On: One-Week Digital Health Challenge}
\begin{enumerate}
  \item Today: Open Screen Time / Digital Wellbeing and review last week's usage data.
  \item Tomorrow: Set daily limits for your most time-consuming entertainment apps (recommend 30-60 minutes).
  \item Day after: Schedule Do Not Disturb (22:00-07:00), allowing calls from favorites.
  \item Day 4: No phone for 1 hour before bed; move charger to the other end of the bedroom.
  \item Day 5: Unfollow 5 accounts or bloggers you no longer watch.
  \item Day 6: Try a full day of "information fast" — no news, no short videos.
  \item Day 7: Review this week's Screen Time report and compare with last week.
\end{enumerate}

\end{tcbpractice}

\begin{tcbsummary}{Chapter Summary}
\begin{itemize}
  \item First check your screen time data — understanding the current state is the first step to change
  \item Set limits for entertainment apps, schedule Do Not Disturb — regain control of your attention
  \item The 20-20-20 rule is the simplest and most effective eye care habit
  \item Looking down at your phone puts 5-6x pressure on your neck — raise your phone and stretch regularly
  \item No phone 1 hour before bed, charger out of reach
  \item The antidote to information anxiety is "active filtering," not "getting more"
\end{itemize}

\end{tcbsummary}
